\section{System Architecture}
\label{sec:system_architecture}

Figure~\ref{fig:hw_architecture} presents the top-level system architecture of the RLCV4.
The diagram captures both the hardware connectivity and the principal data flows through the system. \\

\textbf{Power path.}
A single \SI{24}{\volt} \ac{DC} input feeds a \ac{TVS} protection network, after which two independent synchronous buck converters derive the \SI{3.3}{\volt} logic rail and the \SI{12}{\volt} fan rail.
An optional \SI{24}{\volt} output connector branches from the same protected rail via a full-rated trace to supply daisy-chained downstream fixtures; overcurrent protection is delegated to the upstream \ac{PSU}.
No secondary supply is required. \\

\textbf{Control data path.}
Lighting control data enters the system via two independent input channels.
\ac{DMX}512 frames arrive on the \ac{RS485} physical interface and are decoded by the \ac{MCU}'s \ac{UART} peripheral; the same interface supports bidirectional \ac{RDM} communication when the transceiver direction pin is asserted by the firmware.
Art-Net and \ac{sACN} \ac{UDP} packets are received over the integrated \SI{2.4}{\giga\hertz} \ac{WiFi} radio.
All paths converge inside the ESP32-S3-MINI-1-N8, where the firmware arbitrates between them and computes the output pixel data.\\

\textbf{Output path.}
Computed channel values leave the \ac{MCU} via two output stages operating in parallel.
Addressable \ac{LED} strips receive a single-wire serial bitstream generated by the \ac{RMT} peripheral.
Analogue \ac{RGB} strips are driven by three \ac{MOSFET} stages switched by \ac{PWM} signals from the LEDC peripheral.\\

\textbf{Auxiliary subsystems.}
The \ac{MCU} monitors ambient temperature via a DS18B20 sensor on the \ac{OWI} bus and regulates the cooling fan over a \SI{25}{\kilo\hertz} \ac{PWM} output, with stall detection via the tachometer input.
A 128$\times$32 \ac{OLED} display and a rotary encoder provide the local operator interface over \ac{I2C} and \ac{GPIO} respectively.
A \ac{USB}-C connector exposes the \ac{MCU}'s native \ac{USB} port for firmware programming and serial debug; its VBUS line is isolated from the board power network.\\

The firmware partitions all real-time output tasks onto Core~0 and all \ac{UI} tasks onto Core~1 of the dual-core LX7 processor.
This partitioning is described in detail in Section~\ref{sec:software_architecture}.

\begin{figure}[H]
  \centering
  \begin{tcolorbox}[colback=white, colframe=border, boxrule=0.4mm, arc=5px]
  \resizebox{\linewidth}{!}{%
  \begin{tikzpicture}[>=latex, font=\footnotesize,
      blk/.style={draw, very thick, fill=white, minimum width=2.8cm,
                  minimum height=0.9cm, rounded corners=3pt, align=center},
      pwr/.style={draw, thick, fill=gray!12, minimum width=2.8cm,
                  minimum height=0.9cm, rounded corners=3pt, align=center},
      bus/.style={-latex, thick},
      pwrl/.style={-latex, thick}
    ]

    \draw[dashed, thick, rounded corners=6pt, fill=FrostBg]
      (-0.3, 0.3) rectangle (9.8, 5.4);
    \node[anchor=north west, font=\footnotesize\itshape] at (-0.3, 5.4)
      {Power Domain};

    \node[pwr] (vin)  at (1.4, 3.0) {24\,V\\Input};
    \node[pwr] (prot) at (4.5, 3.0) {TVS + $C_\text{bulk}$\\Protection};
    \node[pwr] (b33)  at (8.0, 4.5) {LMR54410\\Buck~3.3\,V};
    \node[pwr] (b12)  at (8.0, 1.5) {LMR54410\\Buck~12\,V};
    \node[pwr] (chain) at (1.4, 1.5) {24\,V Out\\(Daisy-Chain)};

    \draw[pwrl] (vin.east) -- (prot.west);
    \coordinate (junc) at ($(prot.east) + (0.3, 0)$);
    \draw[pwrl,-] (prot.east) -- (junc);
    \fill         (junc) circle (2.5pt);
    \draw[pwrl, rounded corners=4pt] (junc) |- (b33.west);
    \draw[pwrl, rounded corners=4pt] (junc) |- (b12.west);
    \draw[pwrl, rounded corners=4pt] (junc) -- (6.2, 1.5) -- (chain.east);

    \draw[very thick, rounded corners=5pt, fill=white]
      (11.2, 0.5) rectangle (15.6, 5.4);
    \node[font=\bfseries, anchor=north, inner sep=3pt] at (13.4, 5.3)
      {ESP32-S3-MINI-1-N8};
    \draw[thin] (11.2, 4.65) -- (15.6, 4.65);

    \node[font=\scriptsize, align=left, anchor=north west] at (11.4, 4.55)
      {\texttt{GPIO19/20}~—~USB~D+/D-\\[3pt]
       \texttt{SDA\,/\,SCL}~——~I2C\\[3pt]
       \texttt{RX\,/\,TX}~——~UART (DMX/RDM)\\[3pt]
       \texttt{GPIO}~———~OWI\\[3pt]
       \texttt{GPIO}~———~PWM~/~Tach\\[3pt]
       \texttt{GPIO}~———~LED~Data~/~RGB};

    \draw[pwrl] (b33.east) -- (11.2, 4.5)
      node[midway, above, font=\scriptsize] {3.3\,V};

    \node[blk] (oled)     at (18.4,  5.0) {OLED Display\\(SSD1306)};
    \node[blk] (rs485)    at (18.4,  3.7) {RS-485 Driver\\(MAX3485)};
    \node[blk] (dmx)      at (22.2,  3.7) {DMX512~/~RDM\\Connector};
    \node[blk] (usb)      at (18.4,  2.4) {USB-C\\Connector};
    \node[blk] (temp)     at (18.4,  1.1) {Temp.\ Sensor\\(DS18B20)};
    \node[blk] (ledstrip) at (18.4, -0.2) {Addr.\ LED Strip\\(WS2812B)};
    \node[blk] (rgbpwm)   at (18.4, -1.5) {RGB MOSFETs\\(R~/~G~/~B)};
    \node[blk] (fan)      at (13.4, -0.8) {Cooling Fan\\(4-wire,~12\,V)};

    \draw[latex-latex, thick] (rs485.west) -- (15.6, 3.7)
      node[midway, above, font=\scriptsize] {UART};
    \draw[latex-latex, thick] (dmx.west) -- (rs485.east);
    \draw[bus] (15.6, 5.0) -- (oled.west)
      node[midway, above, font=\scriptsize] {I2C};
    \draw[bus] (15.6, 2.4) -- (usb.west)
      node[midway, above, font=\scriptsize] {USB};
    \draw[bus] (temp.west) -- (15.6, 1.1)
      node[midway, above, font=\scriptsize] {OWI};
    \draw[bus, rounded corners=4pt] (15.2, 0.5) |- (ledstrip.west)
      node[near end, above, font=\scriptsize] {LED Data};
    \draw[bus, rounded corners=4pt] (15.0, 0.5) |- (rgbpwm.west)
      node[near end, above, font=\scriptsize] {RGB PWM};
    \draw[bus] (13.4, 0.5) -- (fan.north)
      node[midway, right, font=\scriptsize] {PWM / Tach};
    \draw[pwrl, rounded corners=4pt] (b12.south) |- (fan.west)
      node[pos=0.15, left, font=\scriptsize] {12\,V};

  \end{tikzpicture}%
  }
  \end{tcolorbox}
  \caption{Top-level system architecture of the RLCV4: power paths (grey nodes), data buses, and output stages.}
  \label{fig:hw_architecture}
\end{figure}
